% objections.tex — Section 6: Objections

\section{Objections}
\label{sec:objections}

\subsection{``Level IV is unfalsifiable''}
\label{subsec:obj-unfalsifiable}

Level~IV itself --- the claim that all consistent mathematical
structures exist --- is not directly falsifiable.  Neither is the
negation.  The claim ``only one structure is physical'' is equally
immune to direct test.

What \emph{is} falsifiable is the derivation chain.  If self-modeling
forces quantum mechanics (Steps~3--6), then discovering a classical
self-modeling system with equal predictive power would refute the
framework.  If the non-composability argument gives $\hthree$, then
discovering a consistent non-special Jordan algebra besides $\hthree$
would break the uniqueness claim (though the classification theorem
makes this impossible).  The specific form of $\rho$ within $\hthree$
(Remark~\ref{rem:rho-h3o}) is under development and may yield
additional falsifiable predictions.

The philosophical premise is not falsifiable.  Its consequences are.
This is the same epistemic status as any foundational principle in
physics: the principle of least action is not falsifiable, but
Lagrangian mechanics is.

\subsection{``Genericity is measure-dependent''}
\label{subsec:obj-measure}

The worry: ``generic among what?  With which measure?''  If the
derivation chain relied on typicality arguments --- ``most structures
are quantum'' --- this would be fatal.  It does not.  The chain
consists of \emph{theorems}, not statistical claims.

Self-modeling forces Jordan algebras (van de Wetering).  Local
tomography forces complex type (dimension counting).  Non-composability
forces $\hthree$ (JvNW + Zel'manov).  At no point does the argument
appeal to ``most structures have property $X$.''  The word ``generic''
appeared in early versions of the program (specifically in the
non-commutativity argument~\citep{motzkintaussky1955}) but has been
superseded.  Paper~5's chain is forcing, not typicality.

\subsection{``This is anthropic reasoning''}
\label{subsec:obj-anthropic}

Anthropic reasoning selects from a pre-existing landscape: ``we
observe quantum mechanics because only quantum-mechanical universes
support observers.''  This requires a landscape (where does it come
from?) and a selection mechanism (what does the selecting?).

For Steps~3--6 (QM), the framework \emph{derives}: self-modeling
forces quantum mechanics by theorem.  There is no landscape to select
from.  The algebraic structure is forced, not filtered.  The
distinction is between ``only habitable houses are inhabited''
(anthropic) and ``this is the only house that can be built on this
foundation'' (forcing).

We acknowledge that Gap~C (complexification, Sec.~\ref{subsec:sm})
\emph{does} use anthropic reasoning: ``no chirality $\to$ no chemistry
$\to$ no self-modelers.''  The program is not uniformly non-anthropic.
Steps~3--6 are forcing; Gap~C is selection.  This is registered
honestly.

\subsection{``The gaps are too big''}
\label{subsec:obj-gaps}

Honest assessment: the gaps are real.  Table~\ref{tab:gaps} registers
every open gap in the program.

\begin{table}[ht]
\centering
\caption{Gap register.  Gaps are classified by severity and shared
  status (whether other programs face the same gap).}
\label{tab:gaps}
\begin{tabular}{@{}p{3.6cm}p{1.2cm}p{2.7cm}@{}}
\toprule
Gap & Severity & Status \\
\midrule
Emergent Lorentz invariance & Standard & Shared by all lattice QG \\
Lattice Bisognano-Wichmann & Standard & Shared by all lattice QG \\
Local equilibrium & Standard & Shared by all lattice QG \\
Continuum limit & Standard & Shared by all lattice QG \\
Gap~C (complexification) & Real & Evolutionary argument; not a thm \\
Minimal composite assumption & Moderate & Load-bearing; not derived \\
Finite-dimensional assumption & Moderate & Natural for subsystems \\
\bottomrule
\end{tabular}
\end{table}

The four GR gaps are standard: every program that derives gravity from
entanglement on a lattice faces them.  They are not bespoke weaknesses
of radical relativity.

Gap~C is the one genuinely novel gap.  The complexification
$\Spin{9} \to \Spin{10}$ is argued on physical grounds (the
observer's $C^*$-nature) but not proved.  What would close it: a
theorem of the form ``$C^*$-measurement maps on a real Jordan module
necessarily factor through $\otimes_{\mathbb{R}} \mathbb{C}$.''  In
the absence of such a theorem, the evolutionary argument (``no
chirality $\to$ no chemistry $\to$ no self-modelers'') provides
overwhelming plausibility without mathematical certainty.

The minimal composite assumption does significant work in the local
tomography proof (Step~4).  All four adversarial reviewers of Paper~5
flagged this.  It is an identified cost, not a hidden one.

\subsection{``How is this different from Tegmark?''}
\label{subsec:obj-tegmark}

Table~\ref{tab:comparison} compares the present framework against
related programs.

\begin{table*}[t]
\centering
\caption{Comparison with related programs.  ``Derives QM'' = quantum
  mechanics follows from the framework's premises.  ``Derives SM'' =
  Standard Model gauge group follows.  ``Measure'' = the framework
  provides a measure on observer-instances.  ``Honest gaps'' = open
  problems are explicitly registered.}
\label{tab:comparison}
\small
\begin{tabular}{@{}p{3.1cm}cccp{2.8cm}@{}}
\toprule
Program & QM & SM & Gaps & Key difference \\
\midrule
\textbf{Radical Rel.} & \textbf{Thm} & \textbf{Cond.}
  & Yes & Forcing, not selection \\
Tegmark MUH & No & No
  & Part. & Elevates math to physical \\
Mueller Alg.\ Idealism & No & No
  & Yes & Measure only, no algebra \\
Wolfram Ruliad & Plaus. & Part.
  & No & ``Can'' not ``must'' \\
String Theory & Input & Landsc.
  & Part. & QM assumed, not derived \\
Barandes Stoch.-Q & Equiv. & No
  & Yes & Why unitary? (we answer) \\
Connes NCG & Input & Cond.
  & Yes & $A_F$ chosen, not derived \\
Adler Trace Dyn. & Emergent & No
  & Yes & Assumes NC; we derive it \\
\bottomrule
\end{tabular}
\end{table*}

\textbf{Tegmark}: elevates mathematics to the status of the physical
(``the physical world IS a mathematical structure'').  We deflate the
physical to the mathematical.  The direction matters: Tegmark inherits
a measure problem (which structure is ``the'' physical one?); we do
not.

\textbf{Mueller}: provides a Solomonoff-type measure on observer
experiences but does not derive the algebraic structure of physics.
Our framework is compatible with Mueller's measure and would benefit
from it as an input to structure space, but the derivation chain does
not require it.

\textbf{Wolfram}: explores all possible computations (the ``ruliad'')
and claims that observers with certain properties will experience
physics resembling ours.  The claim is ``can emerge,'' not ``must
be.''  Our chain is forcing: self-modeling \emph{must} produce complex
QM.

\textbf{Barandes}: establishes a bijection between unitary QM and
indivisible stochastic processes.  This is a load-bearing link in our
chain (Step~7) but does not itself explain \emph{why} unitary dynamics
hold.  Our chain provides the ``why'': self-modeling forces
$C^*$-algebras, whose automorphisms are unitary.

\textbf{Connes}: derives the Standard Model Lagrangian from a spectral
triple, but the finite algebra $A_F = \mathbb{C} \oplus \mathbb{H}
\oplus M_3(\mathbb{C})$ is chosen to match observation.  Our framework
derives $\hthree$ from non-composability and extracts the SM gauge
group as the intersection of two stabilizers.  The algebra is forced,
not chosen.
